\documentclass[11pt]{article}

\usepackage{sectsty}
\usepackage{graphicx}

% Results generated from results/summary.json by the pipeline
\makeatletter%
\newcommand\results[1][all]{%
  \ifnum\pdfstrcmp{#1}{all}=0%
    \def\results@out{%
    \{%
      "mean": {-}0.057822632418548854, %
      "std": 0.9940726228753299%
    \}}%
  \else%
    \ifnum\pdfstrcmp{#1}{mean}=0%
      \def\results@out{%
    {-}0.057822632418548854}%
    \else%
      \ifnum\pdfstrcmp{#1}{std}=0%
      \def\results@out{%
    0.9940726228753299}%
    \else%
      \def\results@out{%
    ??}%
    \fi\fi
  \fi
  \results@out
}%
\makeatother%

% Margins
\topmargin=-0.45in
\evensidemargin=0in
\oddsidemargin=0in
\textwidth=6.5in
\textheight=9.0in
\headsep=0.25in

\title{This is the paper}
\author{ Author }
\date{\today}

\begin{document}
\maketitle
\pagebreak

% Optional TOC
% \tableofcontents
% \pagebreak

%--Paper--

\section{Section 1}

In Figure~\ref{fig:x-vs-y},
we can see the measured relationship between $x$ and $y$,
along with linear and quadratic fits
obtained via ordinary least squares~\cite{draper1998applied}.
The quadratic fit has a much higher $R^2$
(\results[r_squared_quadratic])
than the linear fit
(\results[r_squared_linear]),
indicating that $y$ increases quadratically with $x$.

\begin{figure}[h]
  \centering
  \includegraphics{figures/x-vs-y.png}
  \caption{
    Measured $x$ vs $y$ data with linear and quadratic fits.
    The quadratic fit ($R^2 = \results[r_squared_quadratic]$)
    captures the data much better than the linear fit
    ($R^2 = \results[r_squared_linear]$),
    showing that $y$ increases quadratically with $x$.
  }
  \label{fig:x-vs-y}
\end{figure}


\pagebreak
\section{Section 2}

This is the second section of the paper.

This is another sentence in the second section.

This was added in a PR.

%--/Paper--

\bibliographystyle{plain}
\bibliography{references}

\end{document}
